\begin{abstract}
Global AI weather models such as ECMWF's AIFS now match or beat physics-based forecasts on grid averages, and
zero-shot time-series foundation models such as Google's TimesFM and Amazon's Chronos forecast any series without
training. Neither is built for what many users in India need: the next 24 hours of temperature, humidity and wind
at a specific station. We study whether a small model trained only on a station network's own history can do
better. \pwm is a graph world model of about one million parameters in which every station is a token. It comes in two forms:
a pure form that uses observations only, and a hybrid that also reads global forecasts published before the
forecast starts. We test it in three north-Indian leagues: six airports (METAR), 20 NOAA stations around Delhi
(WMO) and 25 Delhi air-quality stations' weather sensors (CPCB). The opponents are real archived operational
forecasts of AIFS, IFS, ICON, GFS and NOAA's GraphCast-based AIGFS, scored raw and after a trailing local bias
correction, together with TimesFM 2.5 and 3.0 and Chronos-2. Every change from version 0.3 to 0.8 was adopted only
if it beat the default setting on validation by more than two seed standard deviations, and several proposed
changes failed that test. At the airports, the hybrid beats every global model on every variable it forecasts when each is scored on its newest
run published before the forecast starts (AIFS by 4.4\%, 11.8\% and 7.4\%). Against AIFS's freshest run it still
loses on temperature and humidity. The pure model beats all three foundation models on every variable at the
airport and WMO stations. On the noisier CPCB sensors, cross-network training, a 14-day encoder and weight averaging
turned four of its losses into ties, but it still trails Chronos-2 on temperature and all three foundation models on wind. Bias-corrected FourCastNet and GenCast
Mini beat the pure model on airport temperature. We report every result, including the losses, and the test-reuse
risk of a version-by-version study.
\end{abstract}
